\chapter{Introdução e Motivações}
A previsão de séries temporais tem um papel fundamental em diversas áreas, servindo como base para a tomada de decisões, gerenciamento de recursos e planejamento. As aplicações abrangem diversos setores, tais como consumo de energia, diagnósticos médicos, finanças, além da agricultura e da produção industrial \cite{hallSurveyMachineLearning2025, liuForecastMethodsTime2021}. Um dos destaques da área de \textit{Machine Learning} é a técnica de \textit{ensemble learning}, que é capaz de capturar dependências complexas e já é utilizada em inumeras aplicações \cite{carvalhoDynamicSelectionEnsemblebased2025}.

O processo de construção de um \textit{Ensemble}
se dá a partir de três etapas: geração, seleção e combinação \cite{mendes-moreiraEnsembleApproachesRegression2012}. Na primeira etapa, a diversidade entre os modelos base é introduzida por estratégias como a variação das representações de entrada, das inicializações e dos algoritmos de aprendizado, além da amostragem de dados e de atributos \cite{fanDiverseModelsUnited2025}. Essa diversidade atua como uma dimensão complementar do ajuste do modelo, permitindo que os modelos base reduzam erros correlacionados na predição final \cite{woodUnifiedTheoryDiversity}. Na etapa de seleção, o objetivo é filtrar o subconjunto de modelos mais competentes para a predição \cite{koDynamicClassifierSelection2008}. Por fim, na fase de agregação, as previsões dos modelos selecionados são combinadas por meio da média, mediana, votação ponderada, entre outros, com o objetivo de agregar as forças dos preditores \cite{allendeEnsembleMethodsTime2016}.

Apesar de a diversidade ter uma importância reconhecida como uma das peças-chave para o desempenho de um \textit{ensemble}, a literatura ainda não tem um consenso formal sobre como gerênciar e quantificar essa característica do sistema de maneira ideal, abordando-a por meio de heurísticas e métricas isoladas que não demonstraram ter mais sucesso do que outras. Muitos estudos concentram-se em tentar maximizar a variação entre os preditores sem distinção, negligenciando o \textit{trade-off} entre viés, variância e diversidade, e deixando de lado investigações sobre como estruturar a geração e seleção dos modelos base de forma sinérgica \cite{woodUnifiedTheoryDiversity}.

Diante destas limitações, este trabalho propõe uma nova abordagem para a seleção dinâmica de preditores, baseada na modelagem da dinâmica temporal da diversidade e em seu uso antecipativo para equilibrar o \textit{trade-off} entre diversidade e acurácia. Os métodos dinâmicos existentes adaptam pesos ou a composição do ensemble em tempo de execução, seja pelo desempenho recente dos modelos \cite{cerqueira2017}, pelo meta-aprendizado da competência de cada preditor \cite{cerqueira2019}, ou por mecanismos reativos à detecção de drift \cite{saadallah2019}. Nesses casos, a diversidade é, de maneira geral, tratada como uma propriedade avaliada no instante da agregação ou após a mudança já ter ocorrido. No lugar de maximizar a variação entre os preditores ou apenas reagir a ela, a proposta busca explorar a estrutura temporal da diversidade e sua relação com o desempenho do ensemble, de modo a antecipar os momentos em que ela, de fato, contribui para reduzir o erro.

Para guiar este trabalho, as seguintes perguntas de pesquisa foram formuladas:
\begin{itemize}
    \item \textbf{QP1:} A diversidade entre os preditores apresenta uma dependência temporal significativa, e como isso varia conforme a métrica de diversidade e as características da série?
    \item \textbf{QP2:} É possível prever a diversidade futura do ensemble com desempenho superior a uma referência ingênua?
    % A DIVERSIDADE IDEAL É O VALOR DA MÉTRICA DE DIVERSIDADE EM QUE A MÉTRICA DE ERRO UTILIZADA TEM SEU MENOR
    % VALOR
    \item \textbf{QP3:} A distância entre a diversidade do ensemble e a diversidade ideal explica a perda de desempenho, e essa relação depende do regime da série?
    \item \textbf{QP4:} A seleção dinâmica orientada à diversidade ideal estimada supera os métodos dinâmicos de referência?
\end{itemize}

Esta pesquisa busca entregar contribuições significativas para a literatura em três aspectos. O trabalho busca inovar na previsão de séries temporais ao formalizar a diversidade como um processo temporal e preditivo, propondo um mecanismo de seleção de \textit{ensembles} antecipativo. A pesquisa também fornece uma avaliação dessa abordagem frente a mudanças de regime e alta volatilidade, contrastando-a com os principais métodos reativos do estado da arte \cite{cerqueira2019arbitrated}. Por fim, espera-se consolidar a metodologia desenvolvida em código aberto, viabilizando sua aplicação e reprodutibilidade por pesquisadores e profissionais de diversos domínios.

\textbf{Limite: 2 páginas}









